\documentclass[12pt]{article}

% === Packages ===
\usepackage[margin=1in]{geometry}
\usepackage{amsmath, amssymb}
\usepackage{setspace}
\usepackage{titlesec}
\usepackage{enumitem}
\usepackage{helvet}
\renewcommand{\familydefault}{\sfdefault} % modern sans-serif look

% === Formatting tweaks ===
\setstretch{1.15}
\setlist[enumerate]{itemsep=0.4em, topsep=0.2em}
\titleformat{\section}{\Large\bfseries\sffamily}{\thesection.}{1em}{}
\titleformat{\subsection}{\large\bfseries\sffamily}{\thesubsection}{1em}{}

% === Title ===
\title{\vspace{-2em}\textbf{Applied Mathematics Oral Exam Questions and Model Answers}\\
\large Tensors in Fluid--Structure Interaction}
\author{Prepared for Graduate Oral Examination}
\date{}

\begin{document}


\section*{Sample Oral Exam Questions on Tensors}

\subsection*{1. What is the material derivative, and why is it important\\in fluid--structure interaction?}

\textbf{Answer.}
The material derivative measures the rate of change of a quantity as experienced
by a moving fluid parcel. For a scalar field $\phi(x,t)$,
\[
\frac{D\phi}{Dt}
= \frac{\partial \phi}{\partial t} + (u\cdot\nabla)\phi,
\]
and for a vector field,
\[
\frac{D\mathbf{u}}{Dt}
= \frac{\partial \mathbf{u}}{\partial t} + (u\cdot\nabla) \mathbf{u}.
\]
This combines local (Eulerian) time dependence with convective transport.

In IBFE, the material derivative appears in:
\begin{itemize}
    \item the \emph{fluid} momentum equation (left-hand side of Navier--Stokes),
    \item the evolution of the solid configuration along material trajectories,
    \item coupling the Eulerian fluid velocity to the Lagrangian mesh velocity.
\end{itemize}

For a no-slip immersed material, $\partial_t\boldsymbol\chi(X,t)=\mathbf u(\boldsymbol\chi(X,t),t)$; this is a Lagrangian time derivative of the motion, whereas $D\mathbf u/Dt$ is the acceleration field evaluated along a fluid trajectory.

\subsection*{2. Explain the Cauchy stress tensor and how traction arises from it.}

\textbf{Answer.}
The Cauchy stress tensor $\boldsymbol{\sigma}$ describes internal surface
forces per unit current area in a fluid or solid.
For a surface with outward normal $\mathbf{n}$, the traction vector is
\[
\mathbf{t} = \boldsymbol{\sigma}\mathbf{n}.
\]
Each component $\sigma_{ij}$ is the $i$-component of the force acting on a face
whose outward normal is in the $j$-direction.

In IBFE:
\begin{itemize}
    \item fluid and elastic contributions combine in the Eulerian momentum balance,
    \item structural constitutive laws are commonly written with a reference-area stress $\mathbf P$,
    \item the Piola relation $\mathbf P=J\boldsymbol\sigma^{e}\mathbf F^{-T}$ connects elastic Cauchy and first Piola--Kirchhoff stresses.
\end{itemize}

\subsection*{3. Derive the form of the Navier--Stokes equations used in IBFE.}

\textbf{Answer.}
Starting from Cauchy’s equation of motion,
\[
\rho\frac{D\mathbf{u}}{Dt} = \nabla\cdot\boldsymbol{\sigma} + \rho\mathbf{g},
\]
split the stress into pressure and viscous components:
\[
\boldsymbol{\sigma} = -p\mathbf{I} + \boldsymbol{\tau}.
\]
For a Newtonian, incompressible fluid,
\[
\boldsymbol{\tau} = 2\mu \mathbf{D},
\qquad
\mathbf{D} = \tfrac{1}{2}(\nabla \mathbf{u} + \nabla \mathbf{u}^\top).
\]
When $\mu$ is constant, $\nabla\cdot(2\mu\mathbf{D})=\mu\nabla^2\mathbf{u}$ (because $\nabla\cdot(2\mu\mathbf{D}) = \mu\nabla^2\mathbf{u} + \mu\nabla(\nabla\cdot\mathbf{u})$ and $\nabla\cdot\mathbf{u}=0$), and this yields
\[
\rho\left( \frac{\partial \mathbf{u}}{\partial t}
+ (\mathbf{u}\cdot\nabla)\mathbf{u} \right)
= -\nabla p + \mu \nabla^2 \mathbf{u} + \rho \mathbf{g},
\qquad 
\nabla\cdot \mathbf{u} = 0.
\]
The coupled IBFE momentum equation adds an Eulerian elastic force density $\mathbf f^e$ (and any other body force):
\[
\rho\frac{D\mathbf u}{Dt}=-\nabla p+\mu\nabla^2\mathbf u+\mathbf f^e+\rho\mathbf g.
\]
The simple one-field formulation assumes compatible treatment of density and background viscosity. General material contrasts require extensions.

\subsection*{4. How does IBFE represent the deformation of an immersed structure?}

\textbf{Answer.}
The solid is described in a Lagrangian coordinate system $X$ with mapping
$\chi(X,t)$ giving its current position.  
The deformation gradient is
\[
\mathbf{F}(X,t) = \nabla_X \chi(X,t).
\]
The solid finite-element formulation supplies a constitutive law. Examples include hyperelastic, neo-Hookean, and fiber-reinforced materials. The law determines the first Piola--Kirchhoff stress $\mathbf{P}(X,t)$.

In IBFE:
\begin{itemize}
    \item Lagrangian FE computes internal elastic forces,
    \item these forces are \emph{spread} to the Eulerian fluid grid via IB
          delta kernels,
    \item the fluid velocity is interpolated back to the Lagrangian mesh to
          enforce no-slip (kinematic coupling).
\end{itemize}

\subsection*{5. Explain how the force-spreading step works in IBFE.}

\textbf{Answer.}
Given a Lagrangian force density $\boldsymbol{\mathcal F}(X,t)$ per unit reference measure,
the Eulerian force density is
\[
\mathbf{f}(x,t) = \int_{B_0}
\boldsymbol{\mathcal F}(X,t)\,\delta(x - \chi(X,t))\,dX.
\]
Here $B_0$ is the reference (Lagrangian) configuration of the structure, the domain of the material coordinates $X$, and $\chi(X,t)$ is the current position of material point $X$.
In IBFE, $\boldsymbol{\mathcal F}$ is normally defined by a finite-element weak form rather than by pointwise differentiation of a discontinuous stress. Code evaluates the corresponding interaction integrals at quadrature points and replaces the Dirac delta by a regularized IB kernel. Conservation or power consistency follows only when spreading and velocity restriction use compatible weights and adjoint operators.

This step couples the solid’s internal stresses to the fluid’s momentum
equation.

\subsection*{6. What does the no-slip condition mean in IBFE?}

\textbf{Answer.}
IBFE enforces
\[
\frac{\partial \chi}{\partial t}(X,t) = \mathbf{u}(\chi(X,t),t),
\]
i.e.,
the solid moves with the local fluid velocity.
This is implemented by \emph{interpolation} from Eulerian velocity to the
Lagrangian mesh via the same IB kernel used for force spreading.

The immersed interface does not have to conform to the Eulerian grid; its kinematic condition is imposed through interpolation and force coupling. The outer Eulerian domain still requires its own physical boundary conditions (or periodic conditions), and constraints such as contact or prescribed solid motion require additional modeling.

\subsection*{7. What conditions ensure that a matrix of components represents a tensor?}

\textbf{Answer.}
A $3\times 3$ component array $A_{ij}$ represents a second-order tensor if it obeys the appropriate transformation law under a change of orthonormal basis. If the same orthogonal basis change $Q$ is used in the domain and codomain,
\[
A^{(m)} = Q A^{(e)} Q^T.
\]
If it does not transform covariantly with the bases, it is merely an array of numbers tied to a coordinate choice. For a two-point tensor such as the deformation gradient, reference and current bases may change independently, giving $F^{\rm new}=Q_{\rm current}F^{\rm old}Q_{\rm reference}^{T}$.
This is essential in IBFE because deformation gradients, stresses, and rate-of-strain
tensors must be genuine geometric objects independent of the coordinate
representation.

\subsection*{8. Why must the viscous stress tensor be symmetric?}

\textbf{Answer.}
For a classical continuum without body couples or couple stresses, balance of angular momentum implies symmetry of the total Cauchy stress:
\[
\sigma_{ij} = \sigma_{ji}.
\]
For a Newtonian fluid, the viscous Cauchy stress is therefore symmetric as well. (Explicitly, $\boldsymbol\tau=\boldsymbol\sigma+p\mathbf I$ is the difference of two symmetric tensors; and the constitutive law $\boldsymbol\tau=2\mu\mathbf D$ is symmetric by construction, since $\mathbf D$ is the symmetric part of $\nabla\mathbf u$. The antisymmetric part of $\nabla\mathbf u$, the spin, describes local rigid rotation and produces no viscous stress.)

In IBFE, the elastic Cauchy stress is symmetric for a classical non-polar material, but the first Piola--Kirchhoff stress $\mathbf P$ is generally not symmetric because it maps reference normals to current forces. Angular-momentum checks must therefore be made using the correct stress measure and configuration. (Equivalently: $\mathbf P\mathbf F^T = J\boldsymbol\sigma$ is symmetric even though $\mathbf P$ is not.)

\subsection*{References for these answers}
\begin{itemize}
    \item A. F. Bower, \emph{Applied Mechanics of Solids}, CRC Press, 2010 (solidmechanics.org): Appendix B (tensors), \S2.2--2.3 (stress, symmetry).
    \item M. E. Gurtin, E. Fried, L. Anand, \emph{The Mechanics and Thermodynamics of Continua}, Cambridge University Press, 2010.
    \item B. E. Griffith, X. Luo, Hybrid finite difference/finite element immersed boundary method, \emph{Int.\ J.\ Numer.\ Methods Biomed.\ Eng.} 33 (2017) e2888.
    \item C. S. Peskin, The immersed boundary method, \emph{Acta Numerica} 11 (2002) 479--517.
\end{itemize}

\end{document}
